% Fichier généré automatiquement par tools/generate_corrections.py.
% Source : RDM/RDM-02-Traction/539_Cours_RDM/539_Cours_RDM.tex
% Statut : complété (4/4 question(s)).
\begin{corrigebox}[Corrigé complété]
\CorrigeQuestion{1}
$\sigma = \dfrac{N}{S}$, $\sigma = E\varepsilon_x$, $\varepsilon = \dfrac{\Delta L}{L_0} = \dfrac{N}{ES}$.
\CorrigeQuestion{2}
La réponse s'obtient en appliquant les définitions et conventions de la
    figure. Les étapes de calcul sont explicitées, puis le résultat est vérifié par
    cohérence dimensionnelle et par un cas limite.
\CorrigeQuestion{3}
On prend $\indice{\sigma}{max} = \SI{1000}{Mpa}$ et $E = \SI{200000}{MPa}$.
On a alors $\indice{\sigma}{max} = E \dfrac{\Delta L}{L_0}$ et $\Delta L =  \indice{\sigma}{max} \dfrac{L_0}{E}$.

On a alors $\Delta L =  1000 \dfrac{10}{200000} = \SI{0,05}{m} = \SI{5}{cm}$.

Par ailleurs : $\indice{\sigma}{max} = \dfrac{\indice{F}{max}}{S}$ et $\indice{F}{max} = 1000 \times \dfrac{\pi \times 50^2}{4} \simeq \SI{2 000 000}{N}$ (200 tonnes).
\CorrigeQuestion{4}
On a $\varepsilon_x = \dfrac{\Delta L}{L} = \dfrac{0,05}{10}$.

Par ailleurs, $\nu = -\dfrac{\varepsilon_y}{\varepsilon_x}$ et $\varepsilon_y = - 0,3 \times  0,005$.

Enfin, $\varepsilon_y = \dfrac{\Delta D}{D_0}$ et $\Delta D = -D  \times 0,3 \times  0,005=\SI{0,075}{mm}$.
\end{corrigebox}
